% =============================================================================
% Methodology and Experimental Setup — the predictivity sweep
%
% Sources (repo, as of 2026-09-25):
%   plan/small-to-large-predictivity-training-plan.md   design, budgets, data, validation
%   plan/methodology-review-2026-08-21.md               LR law, shallow ladder, checkpoint grid
%   plan/benchmark_selection.md                         benchmark families + per-cell task counts
%   plan/90M-rung-anomaly.md                            the 90M divergence (appendix draft)
%   plan/1b-models.md                                   the 20-checkpoint 1B cells
%   src/pretrain/launch_trainings.py                    DATA_SCHEMES, SEED_TRIPLES, n_checkpoints
%   src/pretrain/hyperparams/hyperparams_{deep,shallow}.json
%   src/pretrain/megatron_args.sh                       every training argument
%   src/pretrain/data/{create_data_mixture,build_data_mixtures}.py, language_sets_scheme*.json
%   src/evals/scripts/score_bpb.py                      BPB
%   src/evals/README.md, src/pretrain/auto_evals_cscs.py harness settings
%   configs/tasks.json                                  the `auto` group and family metadata
%
% Included in the appendix through app_languages.tex. Citation keys follow the
% .bib's Zotero style; keys present in the other sections are reused verbatim,
% the following are GUESSED and must be checked against the .bib:
%   ainslie_gqa_2023 hoffmann_training_2022 bhagia_establishing_2024
%   pagliardini_ademamix_2024 hagele_scaling_2024 gao_framework_2024
%   kwon_efficient_2023 magnusson_datadecide_2025
% Anything still marked \citeneeded{...} has no key yet.
% =============================================================================
\providecommand{\citeneeded}[1]{\textcolor{red}{[CITE: #1]}}
\providecommand{\checknum}[1]{\textcolor{orange}{#1}}


\section{Model Ladder Details}
\label{app:methodology}

We give the training grid and checkpoint schedule (\S\ref{app:training-grid}), model configurations (\S\ref{app:model-configurations}), data construction (\S\ref{app:data-construction}), and evaluation protocol (\S\ref{app:evaluation-protocol}).


\subsection{Training Grid and Checkpoints}
\label{app:training-grid}

\begin{table}[t]
\centering
\small
\setlength{\tabcolsep}{4pt}
\caption{The data schemes of the sweep. Every scheme trains the six-rung ladder from 90M to 1.7B at the settings it defines; the 3B rung is trained only under schemes A and B at $L \in \{8, 15\}$, deep, and is counted in the A and B rows. Seeds: ``grid'' means one seed per cell, or three seeds for the deep cells marked in Table~\ref{tab:grid}; every other scheme trains a single seed. Runs count both architectures where the scheme trains both. A further scheme, BT3 (the scheme-B $L=30$ list at $T=3$), is registered and built but not trained.}
\label{tab:schemes}
\begin{tabular}{llccllr}
\toprule
Scheme & Language lists & $T$ & $L$ & Sizes & Architectures & Runs \\
\midrule
A   & resource-ranked (baseline) & 1 & 1, 2, 8, 15, 30, 50 & all (+3B) & deep, shallow & 92 \\
AT3 & scheme A lists, flattened  & 3 & 15, 30, 50          & all & deep, shallow ($L=50$); deep ($L=15, 30$) & 24 \\
B   & diversity-first            & 1 & 8, 15, 30           & all (+3B) & deep, shallow & 40 \\
ZH  & Chinese as the second language & 1 & 2              & all & deep     & 6 \\
ES  & Spanish as the second language & 1 & 2              & all & deep     & 6 \\
\midrule
Total & & & & & & 168 \\
\bottomrule
\end{tabular}
\end{table}

\textbf{Seeds.} Every cell trains with seed 1904. To measure how much a ranking can move for reasons unrelated to the intervention, a subset of cells trains three seeds, each with a different initialization and data order (Table~\ref{tab:grid}): seeds 64, 313, and 1904 at 175M and 600M for $L \in \{1, 2, 50\}$, and seeds 28, 1797, and 1904 at 1B for $L \in \{1, 2, 30\}$, all in the deep architecture. The seed triples cover 175M, 600M, and 1B; they do not directly estimate seed variability at 1.7B. The extra schemes (AT3, ZH, ES) are intervention levels rather than part of the noise estimate and train a single seed.

\begin{table}[t]
\centering
\small
\caption{The scheme-A deep grid: seeds per $(N, L)$ cell. The shallow architecture trains one seed per cell. The 3B rung (not shown) trains $L \in \{8, 15\}$ with one seed. In total the scheme-A deep grid has 56 runs, the 90M and 3B rungs included.}
\label{tab:grid}
\begin{tabular}{lcccccc}
\toprule
$L$ & 90M & 175M & 350M & 600M & 1B & 1.7B \\
\midrule
1   & 1 & 3 & 1 & 3 & 3 & 1 \\
2   & 1 & 3 & 1 & 3 & 3 & 1 \\
8   & 1 & 1 & 1 & 1 & 1 & 1 \\
15  & 1 & 1 & 1 & 1 & 1 & 1 \\
30  & 1 & 1 & 1 & 1 & 3 & 1 \\
50  & 1 & 3 & 1 & 3 & 1 & 1 \\
\bottomrule
\end{tabular}
\end{table}

\textbf{Checkpoints.} Each run saves evenly spaced checkpoints: 20 per run, 40 at the 1B rung, and 60 at the 1.7B rung, so that the two reference rungs are sampled more densely. Because 40 and 60 are multiples of 20, the denser grids contain the shared fractions $k/20$, for $k=1,\ldots,20$. The $20N$ Chinchilla token budget falls at one fifth of training. We evaluate twelve checkpoints per run: the ten tenths of training, which every size shares and which let us compare sizes at the same multiple of Chinchilla (20\% of a run is 1C and its end 5C), plus the 85\% and 95\% saves. The analysis reads the ten tenths. The two extra points exist so that checkpoint-to-checkpoint noise can be estimated over five evenly spaced checkpoints in the last 20\% of training (80, 85, 90, 95 and 100\%), the decay phase of the schedule.\footnote{Fifteen 1B cells were trained before the 40-checkpoint rule and save 20 checkpoints (every 2,287 steps to step 45,740 rather than every 1,143 steps to 45,720). Their saves land on the same $k/20$ fractions, so they are evaluated at the same points as every other run; the token count at matched checkpoints differs by 0.044\%.}

\textbf{Compute axis.} Every checkpoint has a known $(N, D)$, so results can also be read against training compute. We estimate $\mathrm{FLOPs} = 6\,(N + d_{\mathrm{model}} V)\,D$, where $V$ is the vocabulary size: we omit embedding-lookup cost but include the output projection, and with a 131K vocabulary the projection nearly doubles the per-token cost of the smallest rungs.

\subsection{Model Configurations and Optimization}
\label{app:model-configurations}

\textbf{Parameter-count convention.} All sizes are non-embedding parameters, that is, the transformer blocks alone, excluding the tied token embedding and output projection. This follows the ladder convention of \citep{bhagia_establishing_2024} adopted by \citet{heineman_signal_2025}. The convention matters here because the 131,072-token vocabulary makes the tied embedding matrix $d_{\mathrm{model}} \times V$ parameters on its own, 101M at the 90M rung and 302M at the 1.7B rung; including it would substantially change the relative spacing of the smallest models on the size axis.

\textbf{Architecture.} All models are decoder-only transformers following the Apertus pretraining recipe \citep{apertus_apertus_2025}: grouped-query attention \citep{ainslie_gqa_2023} with four query heads per key-value head and a head dimension of 64, rotary position embeddings \citep{su_roformer_2023} with base 500,000 and Llama-3-style scaling factor 8, RMSNorm, QK-layernorm, the xIELU activation \citep{huang_xielu_2025}, no bias terms, no dropout, a sequence length of 4,096, and tied input and output embeddings over the 131,072-token Apertus tokenizer (\texttt{swiss-ai/Apertus-70B-2509}). The deep ladder scales by jointly increasing depth, width, and the number of attention heads, keeping a width-to-depth ratio $d_{\mathrm{model}}/n_{\mathrm{layers}}$ of 64 from 175M to 1B (51 at 90M and 77 at 1.7B, the endpoint compromises needed to hit the target sizes) and a feed-forward multiplier of 4. Table~\ref{tab:ladder} gives the configurations.

\textbf{Shallow ladder.} The depth intervention trains a second ladder at the same six non-embedding sizes with roughly twice the width-to-depth ratio (110 to 146 against 51 to 77). To constrain the architecture comparison, the shallow search pins the feed-forward multiplier and the query-to-key-value head ratio to the deep values, restricts $d_{\mathrm{model}}$ to multiples of 256, and among candidates within 4\% of the target size picks the ratio closest to 128. The shallow sizes land within 0.4\% to 5.2\% of the deep ones, and each shallow run trains its own $D = 100N$.

\begin{table}[t]
\centering
\small
\setlength{\tabcolsep}{3.5pt}
\caption{The model ladder. $N$ counts non-embedding parameters; the tied embedding adds $d_{\mathrm{model}} \times 131{,}072$. $D = 100N$ is the token budget; steps at 504 sequences of 4,096 tokens (2.06M tokens per step); LR is the peak learning rate from the compute-based law at the run's own budget; ckpts is the number of saved checkpoints. GQA (grouped-query attention) gives query/key-value head counts. Shallow rows approximately match the deep models in $N$; their token budgets follow their own parameter counts. FFN is the feed-forward hidden dimension.}
\label{tab:ladder}
\begin{tabular}{llrrrrrrrrr}
\toprule
Arch & Size & Layers & $d_{\mathrm{model}}$ & FFN & GQA & $N$ (M) & $D$ (B) & Steps & LR ($10^{-3}$) & Ckpts \\
\midrule
deep & 90M  & 15 &  768 & 3,072 & 12/3  &   92.9 &   9.3 &  4,500 & 1.43 & 20 \\
deep & 175M & 16 & 1,024 & 4,096 & 16/4  &  176.2 &  17.6 &  8,540 & 1.22 & 20 \\
deep & 350M & 20 & 1,280 & 5,120 & 20/5  &  344.1 &  34.4 & 16,660 & 1.03 & 20 \\
deep & 600M & 24 & 1,536 & 6,144 & 24/6  &  594.5 &  59.5 & 28,800 & 0.90 & 20 \\
deep & 1B   & 28 & 1,792 & 7,168 & 28/7  &  944.1 &  94.4 & 45,720 & 0.80 & 40 \\
deep & 1.7B & 30 & 2,304 & 9,216 & 36/9  & 1,672.2 & 167.2 & 81,000 & 0.69 & 60 \\
\midrule
shallow & 90M  &  8 & 1,024 &  4,096 & 16/4  &   88.1 &   8.8 &  4,260 & 1.45 & 20 \\
shallow & 175M & 10 & 1,280 &  5,120 & 20/5  &  172.0 &  17.2 &  8,340 & 1.22 & 20 \\
shallow & 350M & 14 & 1,536 &  6,144 & 24/6  &  346.8 &  34.7 & 16,800 & 1.03 & 20 \\
shallow & 600M & 14 & 2,048 &  8,192 & 32/8  &  616.6 &  61.7 & 29,860 & 0.89 & 20 \\
shallow & 1B   & 17 & 2,304 &  9,216 & 36/9  &  947.6 &  94.8 & 45,920 & 0.80 & 40 \\
shallow & 1.7B & 20 & 2,816 & 11,264 & 44/11 & 1,665.3 & 166.5 & 80,640 & 0.69 & 60 \\
\bottomrule
\end{tabular}
\end{table}

\textbf{Optimization.} Every run uses the same optimizer and schedule, with only the per-size values changing. We train with AdEMAMix \citep{pagliardini_ademamix_2024} ($\beta_1 = 0.9$, $\beta_2 = 0.999$, $\beta_3 = 0.9999$, $\alpha = 8$), weight decay 0.1, gradient clipping at 0.1, and a global batch of 504 sequences of 4,096 tokens, about 2.06M tokens per step, at every size. The learning rate follows a warmup-stable-decay schedule \citep{hagele_scaling_2024}: linear warmup over about 4\% of the steps, a constant phase, and a $1-\sqrt{\cdot}$ decay to zero over the final 20\%. The peak learning rate comes from a compute-based law evaluated at each run's own budget, $\eta = 0.3118\, C^{-1/8}$ with $C = 6 N D = 600 N^2$, so it decreases smoothly from $1.43 \times 10^{-3}$ at 90M to $6.9 \times 10^{-4}$ at 1.7B. We scale optimizer warmup with run length and initialization with width: the AdEMAMix $\alpha$ and $\beta_3$ warmup spans the full run, and the initialization standard deviation is width-scaled as $0.008944 \sqrt{1792 / d_{\mathrm{model}}}$. Training uses BF16 mixed precision with FP32 gradient accumulation in the Swiss AI fork of Megatron-LM \citep{shoeybi_megatron-lm_2019}, with data parallelism only (no tensor or pipeline parallelism).

\textbf{The 90M rung.} Nine of the ten completed 90M runs diverge: each reaches its best training loss between 15\% and 19\% of the way through training and degrades for the rest of the run, ending 1.2 to 1.9 nats above its own best. No larger rung shows this. Held-out bits-per-byte rises with training on every language we checked, so the degradation is real rather than an artefact of the training objective. Retraining the rung at three learning rates spanning a factor of five and at one half and one sixth of the batch size reproduces the failure at the same fraction of the optimizer steps, which points to the interaction between the fixed $\beta_3$ timescale of AdEMAMix, about 10,000 steps, and a run that is only 4,500 steps long. We keep the rung in the grid and report its results, but exclude it from the scaling fits. Appendix~\ref{app:90m-rung} documents the investigation.


\subsection{Data Construction}
\label{app:data-construction}

\textbf{Language lists.} The language lists are generated from the FineWeb-2 per-language byte counts and the benchmark availability in the evaluation harness. \emph{Scheme A} ranks FineWeb-2 subsets by training-split UTF-8 bytes, excluding the unidentified-language subsets, and takes the top $L-1$; the lists are nested, so the $L=2$ list (Russian) is a subset of the $L=8$ list (Russian, Chinese, German, Japanese, Spanish, French, Italian), which is a subset of the $L=15$, $L=30$, and $L=50$ lists. A 100-language list was generated by the same rule (with the eight subsets among the top 99 that have no benchmark in the harness replaced by the next subsets by bytes with at least two benchmark families) but was not trained, as explained below. \emph{Scheme B} builds diversity-first lists at $L \in \{8, 15, 30\}$ within the top 49: it takes the largest subset of each script not yet covered, then of each language family not yet covered, then the remaining subsets by bytes. At $L=8$ the scheme-B list spans seven scripts against four for scheme A, and at $L=30$ it spans 14 scripts and 12 families against 10 and 10. \emph{Schemes ZH and ES} replace Russian with Chinese or Spanish at $L=2$. The full lists are given in Appendix~\ref{app:languages}.

\textbf{Mixture composition.} Every multilingual setting blends English and FineWeb-2 50/50 at training time with the Megatron data loader's blend weights, from one English dataset and one FineWeb-2 dataset per (scheme, setting). Within the FineWeb-2 half, language $i$ receives a share $q_i \propto p_i^{1/T}$, where $p_i$ is its estimated token proportion in the corpus. Scheme A allocates at $T=1$, proportionally to the corpus. The FineWeb-2 byte distribution is extreme: on the planned 100-language list at $T=1$ the largest language would take 28\% of the multilingual half and the smallest 0.0017\%, and at the 90M rung 66 of the 99 languages would receive under 10M tokens. The flattened scheme AT3 allocates at $T=3$, which would lift the median language of that build from 90M to 373M tokens. The $L=100$ setting was dropped from the grid; the temperature intervention is instead trained on the scheme-A lists at $L=50$ (both architectures) and, where flattening starves no language, at $L=30$ and $L=15$ (deep only), whose smallest language keeps 343M and 1.2B tokens at $T=1$ under the 1.7B budget.

\textbf{Build sizing and repetition.} Each FineWeb-2 dataset is built to half of the largest budget that trains on it plus 10\% headroom, 92B tokens for the 1.7B rung; the ZH build (52B) was sized when that scheme stopped at 1B and is read unchanged by its 1.7B cell. The English dataset is built once to 184B tokens. The builder never repeats a document: a language that runs out of data yields a smaller build. Smaller models draw a shuffled fraction of the same build, so within a setting all sizes see the same language proportions. Two departures from a single pass are recorded and carried into the analysis. First, no single FineWeb-2 subset can feed the 1.7B rung at $L=2$: the Russian build realizes 72.8B tokens against the 83.6B a 1.7B draws from the multilingual half, so the 1.7B at $L=2$ sees Russian about 1.15 times. The Chinese build (52.0B tokens) and the Spanish build (23.9B, all the Spanish the source holds) are smaller still: at 1.7B, Chinese is seen 1.61 times and Spanish 3.51 times (0.65, 0.91 and 1.98 times for Russian, Chinese and Spanish at 1B). Repetition stays below 1.7 epochs across the Russian and Chinese ladders but rises from 0.7 at 350M to 3.5 at 1.7B for Spanish, so any comparison of the $L=2$ schemes states these epoch counts. Second, the $L=15$ and $L=50$ builds of scheme A and the $L=15$ build of scheme B were originally built to 52B and extended to 92B when the 1.7B rung was added at those settings; the extension continues each language's token stream from the same starting point, so the six 1.7B cells there read a superset of what their proxies read, extended with later CommonCrawl dumps.\footnote{Verified byte for byte: every language section of each 52B build is a prefix of the corresponding 92B build.}

\textbf{Validation set.} One fixed validation set is shared by every model and every scheme. For each of the 99 FineWeb-2 languages of the $L=100$ list plus English, it takes the first 5M tokens of that language's first parquet file, capped at 30\% of the file's rows so that single-file languages keep training data. The build records each language's token count, its UTF-8 byte count, and the number of leading rows assigned to validation; every training build reads this manifest and skips exactly those rows, so training and validation are disjoint by construction. Since there is one validation build for the whole sweep, bits-per-byte is comparable across sizes, language settings, and schemes.

\subsection{Evaluation Protocol}
\label{app:evaluation-protocol}

\textbf{Bits-per-byte implementation.} We measure the negative log-likelihood of held-out text in bits, normalized by its UTF-8 byte count.
For language $\ell$, let $x_1,\ldots,x_{m_\ell}$ be the evaluated token sequence, $c_t$ the context preceding token $x_t$ within the scoring block, and $b_\ell$ the number of bytes in the evaluated text. We compute
\begin{equation}
\mathrm{BPB}_\ell = \frac{-\sum_{t=2}^{m_\ell}\log_2 p(x_t \mid c_t)}{b_\ell}.
\label{eq:bpb}
\end{equation}

We concatenate validation documents without end-of-document separators and score them in blocks with 4,096 input tokens and one-token overlap between successive blocks. Each token except the first of the stream is scored once; context resets at block boundaries. We select a deterministic prefix ending at the document boundary that reaches 1M tokens, or the full stream if it is shorter. We measure the denominator by decoding the entire selected prefix and counting its UTF-8 bytes, including the initial conditioning token. Thus, the implementation omits the first token's likelihood but retains its bytes. We also compute token-level perplexity from the same forward passes.

% The original draft reports a standard error of about 0.003 bits for this
% prefix. Add the estimation procedure before presenting it as uncertainty.

\begin{table}[t]
\centering
\small
\setlength{\tabcolsep}{4pt}
\caption{Benchmark families evaluated during pretraining. ``Langs'' is the number of languages the harness ships for the family; a model is evaluated only on the languages in its training mixture. MC = multiple choice; QA = question answering; UD = Universal Dependencies. Construction: HT = human or professional translation, MT = machine translation, native = written in the language, auto = automatically generated.}
\label{tab:benchmark-families}
\begin{tabular}{llllr}
\toprule
Family & Capability & Format & Construction & Langs \\
\midrule
Belebele          & reading comprehension     & 4-way MC, 900 items/lang      & HT (FLORES)           & 122 \\
Global PIQA           & physical commonsense      & 2-way completion              & native                & 116 \\
MultiBLiMP            & grammatical acceptability & minimal pairs                 & auto (UD, UniMorph)   & 101 \\
INCLUDE             & regional knowledge        & MC exam questions             & native                &  44 \\
Global-MMLU  & world knowledge & MC (MMLU)                & HT + community        &  42 \\
Okapi HellaSwag           & activity commonsense      & 4-way completion              & MT                    &  31 \\
Okapi ARC             & science QA                & MC                            & MT                    &  31 \\
XNLI                 & natural language inference & 3-way                        & HT                    &  15 \\
XStoryCloze              & narrative commonsense     & 2-way ending                  & HT                    &  11 \\
XCOPA                  & causal commonsense        & 2-way, 500 items              & HT                    &  11 \\
PAWS-X                 & paraphrase identification & 2-way                         & HT                    &   7 \\
XWinograd     & coreference               & Winograd schemas              & HT                    &   6 \\
LAMBADA multilingual & last-word prediction & completion             & MT                    &   5 \\
TruthfulQA-Multi  & truthfulness           & MC                            & HT                    &   5 \\
\bottomrule
\end{tabular}
\end{table}

\textbf{Task selection.} We select the registered tasks whose language is in the training mixture; task counts can exceed family counts. For the resource-ranked lists, the number of tasks grows with $L$: 106 at $L=1$, 125 at $L=2$, 294 at $L=8$, 430 at $L=15$, 590 at $L=30$, and 776 at $L=50$. These counts cover the whole evaluated group, which besides the multilingual families of Table~\ref{tab:benchmark-families} contains the reformulated variants of the letter-format families (\S\ref{subsec:rq1}) and a set of English-only benchmarks \checknum{(the table predates their addition and must be regenerated)}. Every trained language is covered by at least one family. For the seed-1904 deep scheme-A ladder we additionally run every task in every language, independent of the training mixture, to read cross-lingual transfer on the benchmarks as well as on bits-per-byte. Appendix~\ref{app:benchmark-coverage} gives the per-language coverage.

\textbf{Harness settings.} All benchmarks run through the Swiss AI fork of the LM Evaluation Harness \citep{gao_framework_2024}, installed from a pinned shared checkout with the source commit recorded in each evaluation job, with the vLLM backend \citep{kwon_efficient_2023}, a maximum context of 4,096 tokens, a beginning-of-sequence token prepended to every prompt, no chat template, and each task's default prompt and number of in-context examples. We report accuracy for multiple-choice and completion tasks and exact match for the generative ones, and aggregate sub-tasks (for example the MMLU subjects) to their family by the harness's own grouping. Base models at these sizes are at chance on several families; the above-chance gate of Appendix~\ref{app:results-filtering} decides per task and size whether such a cell enters the analysis.
